\Needspace{34\baselineskip}
\begin{center}
GAUSS TO VON ZACH.
\end{center}

\begin{center}
Monatliche Correspondenz. Vol. XV. Pp. 152–157. February 1807.
\end{center}

\begin{flushright}
Brunswick, 3 January 1807.
\end{flushright}

\par\indent --- --- --- I have the honor of sending you herewith the ephemeris for the course of \emph{Ceres}
for want of other observations\HistoricalNote{26}, I could base the elements for calculating it only on Professor
\SourceLeaf{286}%
\textsc{Pasquich}’s observation of the last opposition (M.\,C., vol.~XIII, \href{https://www.digitale-sammlungen.de/en/view/bsb10538605?page=200}{p.~192}; vol.~XIV, \href{https://www.digitale-sammlungen.de/en/view/bsb10538606?page=102}{p.~92}), so I must put off improving them further until I receive more reliable observations.
They were as follows:

\begin{center}
XI. \emph{Elements.}
\end{center}

\begin{center}\begingroup
\fontsize{9}{11}\selectfont
\setlength{\tabcolsep}{3pt}\renewcommand{\arraystretch}{1.12}
\begin{tabular}{ll}
\multicolumn{2}{l}{Epoch of longitude on the meridian of} \\
\quad Seeberg 1806 .\ .\ .\ .\ .\ .\ .\ . & $108^\circ\,19'\,34''7$ \\
Longitude of aphelion in 1806 .\ .\ . & $326\;37\;59$ \\
Longitude of the ascending node$\cdot$1806  & $80\;53\;23$ \\
Inclination of the orbit .\ .\ .\ .\ .\ .\ . & $10\;37\;33.7$ \\
Mean annual tropical motion  & $78^\circ\,9'\,23''3$ \\
Mean daily tropical motion . & $770''8584$ \\
Eccentricity .\ .\ .\ .\ .\ .\ .\ .\ .\ . & $0.0783486$ \\
Log. of the semiaxis .\ .\ .\ .\ .\ . & $0.4420728$ \\
\end{tabular}\endgroup\end{center}

\par\indent My little work on the determination of planetary orbits and so on, to which I have devoted most of my time, is nearing completion. So far I have written 22 sheets of manuscript, and I am now putting in the detailed examples. One is taken from \emph{Juno}, where the interval is 22 days. In the second, from \emph{Pallas}, it is 71 days. I have since struck out the third, from \emph{Ceres}, where it was 118 days, so that another, where it is 260 days, can bring out the great generality of the method even better. In all these examples the method can be applied directly, with nothing at all about the orbit taken as already known. The whole will amount to about 30 sheets in print. In the present circumstances, I fear it may take some effort to find a reliable publisher in Germany. --- --- --- [The geocentric course of Ceres from 19 December 1806 to 21 September 1807 is shown just as I indicated above for Pallas on \hyperlink{native-283}{p.~278} of this volume, with the sole difference that the brightness is given to five decimal places, where it was given to four.]
